\section{从平行采样走向可回溯推理}

课堂讨论提出两条改进路线：简单题少采样、困难题多采样；以及把“先想算法再写代码”的 reasoning 直接纳入训练与搜索。

对复杂任务，可以把一次完整程序生成改写成分层搜索：

\begin{enumerate}
    \item 生成若干算法草案与复杂度分析；
    \item 对每个草案生成关键不变量和边界测试；
    \item 只为高潜力草案生成完整代码；
    \item 执行失败后定位到相应阶段，允许回溯而非整段重写。
\end{enumerate}

\begin{importantbox}{树搜索需要可恢复的中间状态}
如果每一步只输出不可解释的长文本，系统难以判断应继续、修正还是回退。结构化计划、显式假设、测试结果和工具状态是长轨迹搜索的基本接口。
\end{importantbox}

\begin{warningbox}{分解并不自动解决错误传播}
第一步的错误算法可能让后续所有步骤“局部正确、全局错误”。每个子目标都需要验证，系统还要允许回到上层重新选择分支。
\end{warningbox}

\subsection{本章小结}

大规模平行采样擅长覆盖，分层 reasoning 与回溯擅长利用反馈。下一代 coding agent 会把两者结合：先探索算法族，再沿少量路径执行、验证和修复。

